\begin{figure}[htbp]
\centering
\begin{tikzpicture}[x=.58cm,y=.58cm]
 \foreach \etapa/\dx in {2/0,4/10.4}{
  \begin{scope}[shift={(\dx,0)}]
   \bcuadricula{6}
   \breina{6}{1}{2}
   \breina{6}{2}{4}
   \ifnum\etapa=2\relax
    \def\casillas{3/1,3/6,4/1,4/3,5/3,5/5,6/1,6/3,6/5,6/6}
    \def\corte{b=2/3}
   \else
    \breina{6}{3}{6}
    \breina{6}{4}{1}
    \def\casillas{5/3,5/5,6/5}
    \def\corte{b'=1/3}
   \fi
   \foreach \fila/\col in \casillas{
    \draw[btinta,line width=.55pt] (\col-.5,6.5-\fila) circle[radius=.16];
    \bmarca{6}{\fila}{\col}
   }
   \draw[btinta,line width=.95pt] (0,6-\etapa)--(6,6-\etapa);
   \node[font=\small,anchor=north] at (3,-.18) {$\corte,\quad m=\the\numexpr6-\etapa\relax$};
  \end{scope}
 }
 \draw[->,>=stealth,btinta,line width=.8pt] (6.55,3)--(9.45,3)
   node[midway,above,font=\scriptsize] {filas 3 y 4};
\end{tikzpicture}
\caption{Dos cortes de la solución \((2,4,6,1,3,5)\). La entrada \(r\) da la columna de la reina
en la fila \(r\). En el ejemplo, las filas se revelan de arriba abajo y la raya horizontal marca hasta dónde
se ha llegado. Las coronas indican las reinas reveladas y los círculos con punto, las casillas
que esas reinas dejan disponibles. Al revelar las filas \(3\) y \(4\), quedan tres de las diez.}
\label{fig:cortes}
\end{figure}
